% Unnumbered full candidate, 2026-11-19 after checkpoint 13.
% No promotion or movement before root's complete independent audit.
% AI-produced by GPT-6 Astra (OpenAI, ChatGPT Pro mode), inherited effort.
\paragraph{Candidate: support-localized ordinary-loop realization in dimension four.}
Let $B$ be a connected closed oriented real PL four-manifold.
Let $j:W\hookrightarrow B$ be a compact connected simply connected
codimension-zero PL submanifold, with collared boundary. The boundary
may be empty, in which case $W=B$. Write
\[
 j_!:H^2(W,\partial W;\mathbb Z/2)
                \longrightarrow H^2(B;\mathbb Z/2)
\]
for extension by zero, and write $\mathrm{PD}_W$ for relative
Poincar\'e--Lefschetz duality. For an ordinary based TOP loop
$\gamma$ of automorphisms of $B$, define
\[
 \eta_B:\pi_1 A_{\mathrm{TOP}}(B)\longrightarrow H^2(B;\mathbb Z/2),
 \qquad \eta_B[\gamma]=\kappa_B(T_\gamma),
 \qquad T_\gamma(x,t)=(\gamma_t(x),t).
\]
This is the A.136 trace homomorphism, also displayed in
(A.140.6); its definition does not require an assembly hypothesis.
Retain $I_B,J_B,c_P$ from A.137.
Then
\[
 \operatorname{im}\eta_B\supseteq\operatorname{im}j_!
   =\mathrm{PD}_B^{-1}\!\left(
       \operatorname{im}\bigl(j_*:H_2(W;\mathbb Z/2)
                            \to H_2(B;\mathbb Z/2)\bigr)\right).
                                                               \tag{L.1}
\]
Every class in the displayed subgroup is realized by an ordinary
TOP loop on $B$ that is the identity outside $\operatorname{int}W$
and on one common collar of $\partial W$ throughout the loop.
No assembly hypothesis on $\pi_1(B)$ is imposed.

If $j_*$ on $H_2(-;\mathbb Z/2)$ is onto, then
\[
 \operatorname{im}\eta_B=J_B=H^2(B;\mathbb Z/2),
 \qquad
 \operatorname{im}(\pi_0c_P)=I_B.                    \tag{L.2}
\]
In that case the automorphism comparison
$A_{\mathrm{PL}}(B)\to A_{\mathrm{TOP}}(B)$ fails to be
surjective on $\pi_1$ whenever $H^2(B;\mathbb Z/2)\ne0$.

In particular (L.2) holds for
\[
 B=C\mathbin{\#}D,
\]
where $C$ is a closed connected simply connected oriented PL
four-manifold and $D$ is a closed connected oriented PL
four-manifold with $H_2(D;\mathbb Z/2)=0$. It therefore holds for
$C\mathbin{\#}\,k(S^1\times S^3)$ for every ordinary finite
integer $k\ge0$.
The equality $\operatorname{im}(\pi_0c_P)=I_B$ retains the
endpoint condition: a TOP concordance component is in the PL
image exactly when its top is TOP-isotopic to a PL automorphism.
We do not assert that this endpoint condition is automatic for
these nonsimply connected $B$, and do not compute either whole
concordance space or $\pi_0P_{\mathrm{PL}}(B)$.

\begin{proof}
The case $W=B$ is A.140. We give all relative-boundary details
for the remaining case.

\textbf{The relative invariant on $W$.}
Put
\[
 N_W=W\times I,
 \qquad
 \partial N_W=(W\times\partial I)\cup(\partial W\times I).
\]
Products with corners are understood with their usual compatible
collars; the underlying PL manifold with boundary can equally be
rounded once. Fix its product PL structure $\alpha_W$ induced
by the given structure on $W$. For a TOP homeomorphism $F$ of
$N_W$ that is the identity on all of $\partial N_W$, define
\[
 \kappa_W(F)=[F_*\alpha_W]-[\alpha_W]
       \in H^3(N_W,\partial N_W;\mathbb Z/2)
       \cong H^2(W,\partial W;\mathbb Z/2).           \tag{L.3}
\]
This is the same relative dimension-five PL-lift difference
used in A.137. The original input is Kirby--Siebenmann,
Essay~V, Theorem~3.2 and its relative convention on p.~240,
together with the universal KS square in \S5.6,
pp.~252--253 \cite{KirbySiebenmann1977}. The marked set here
is the whole boundary, including its vertical part. A collar
extension to an open five-manifold gives exactly the source's
relative closed-set formulation. Thus a boundary-fixed isotopy
preserves (L.3).

The cocycle is additive for the maps used here. Indeed, any
bottom-fixed homeomorphism $F$ of $W\times I$ induces the
identity on absolute homology: inclusion of the bottom $W$
is a homotopy equivalence, and $F$ restricts there to the
identity. By naturality of mod-$2$ Poincar\'e--Lefschetz
duality, $F$ therefore acts trivially on
$H^3(N_W,\partial N_W;\mathbb Z/2)$ as well.
The natural affine-lift cocycle, as in A.137, gives
\[
 \kappa_W(FG)=\kappa_W(F)+\kappa_W(G),
 \qquad \kappa_W(F^{-1})=\kappa_W(F).                 \tag{L.4}
\]
This argument is valid with $\partial W\ne\varnothing$;
it does not discard the vertical boundary condition.

\textbf{Every relative class is realized by an ordinary loop.}
Choose one compatible smoothing of $W$. The low-dimensional
existence statement in Kirby--Siebenmann, Essay~I, p.~39,
applies to this compact PL manifold; no isotopy or family is
being smoothed. The induced PL lift is the chosen standard
one in (L.3).
Orson--Powell--Randal-Williams explicitly work with compact
smooth four-manifolds with boundary and maps relative to that
boundary: see \S1.1, p.~1, and \S2.1--2.2, pp.~4--5
\cite{OrsonPowellRandalWilliams2026}.

The comparison of their invariant with (L.3) has the relative
form
\[
       \mathrm{KS}_{\mathrm{OPR}}(F)
          =\mathrm{PD}_W\bigl(\kappa_W(F)\bigr)
                   \in H_2(W;\mathbb Z/2).           \tag{L.5}
\]
For completeness, use $Z=N_W\times I$ and
$C_Z=(N_W\times\{0\})\cup(\partial N_W\times I)$.
The standard lift extends over $Z$. Changing only its top
boundary lift gives the difference class in
$H^3(\partial Z,C_Z)\cong H^3(N_W,\partial N_W)$.
Its obstruction is the triple connecting isomorphism to
$H^4(Z,\partial Z)$, because $C_Z\hookrightarrow Z$
is a homotopy equivalence. This is the interval suspension.
OPR uses the pulled rather than the pushed structure;
(L.4) makes their difference classes agree over $\mathbb Z/2$.
The two interval suspensions, followed by duality on $Z$,
are precisely duality
$H^2(W,\partial W)\to H_2(W)$. This proves (L.5) as an
equality of classes. It is the same universal-lift argument
as A.140, now retaining all vertical boundary faces.

Since $W$ is simply connected, the assembly domain in
OPR's Definition~5.8 is zero, and the Bockstein target
$H_1(W;\mathbb Z_{(2)})$ is zero. Their Theorem~2.3(D),
p.~6, therefore realizes every $z\in H_2(W;\mathbb Z/2)$
by a literal boundary-fixed TOP homeomorphism $F_0$ of
$N_W$. To see the literal boundary condition without
identifying block loops with ordinary loops, repeat the
normalization in A.140: the block realization first gives
$F$ from identity to a diffeomorphism $f$ that is smoothly
pseudo-isotopic to identity relative to $\partial W$.
If $G$ is such a pseudo-isotopy from $f$ to identity, then
$G'=G\circ(f^{-1}\times1)$ is smooth, bottom-fixed,
and has top $f^{-1}$. Thus $F_0=G'F$ fixes the whole
boundary of $N_W$ and has the same OPR class, by source
additivity and the zero obstruction of the smooth $G'$.
Equation (L.5) realizes every desired class
$v\in H^2(W,\partial W;\mathbb Z/2)$ as $\kappa_W(F_0)$.

The corrected Gabai--Gay--Hartman--Krushkal--Powell
Theorem~1.5, p.~2, applies to compact simply connected
topological four-manifolds with boundary
\cite{GabaiGayHartmanKrushkalPowell2026}.
It gives an ordinary path $H_s$ from identity to $F_0$
through bottom-fixed, vertical-boundary-fixed TOP
pseudo-isotopies. Its top restrictions $\gamma_s$
form an ordinary based loop of automorphisms of $W$
fixing $\partial W$ pointwise. As in A.140 put
\[
 T_s(x,t)=(\gamma_{st}(x),t),\qquad Q_s=T_s^{-1}H_s.
\]
Every $Q_s$ fixes the entire boundary of $N_W$.
It joins identity to $T_\gamma^{-1}F_0$.
Apply common-collar normalization to this whole family;
relative lift invariance and (L.4) give
\[
                \kappa_W(T_\gamma)=v.                \tag{L.6}
\]
No ordinary loop is extracted from the block theorem
until this independent ordinary pseudo-isotopy theorem
has supplied the path $H_s$.

\textbf{Common support and extension by identity.}
We may arrange that $\gamma$ is the identity on one
fixed collar of $\partial W$ for all loop parameters.
Here is the elementary family construction, so that this
does not require a dimension-four classification theorem.
Choose a PL collar $\partial W\times[0,d]\hookrightarrow W$.
For $0\le a\le d/4$, define an embedding $c_a:W\to W$
by identity outside that collar and, inside it, by
\[
 (x,t)\longmapsto(x,\sigma_a(t)),\qquad
 \sigma_a(t)=\min\{t+a,(t+d)/2\}.
\]
It maps the collar interval onto $[a,d]$, and $c_0=1$.
Conjugate each $\gamma_s$ by $c_a$ on $c_a(W)$ and
extend by identity on the removed collar of width $a$.
The formulas glue because every $\gamma_s$ fixes
$\partial W$. They give a based loop homotopy, and
preserve the identity loop exactly. At $a=0$ joint
continuity, including inverse continuity, follows from
uniform continuity of the given compact family on the
compact manifold and its fixed boundary. Equivalently,
the same collar compression was verified in A.135--A.136.
At $a=d/4$ the whole loop has one common identity collar.

Reparameterize the loop by a PL function $\rho:I\to I$
equal to zero near $0$ and to one near $1$. This is a
based loop homotopy, using
$(1-u)t+u\rho(t)$, so it preserves (L.6). Its trace is
now identity near every face of $\partial N_W$.
Extend the loop by identity to $B$. The finite closed-set
gluing of $W$ and $B\setminus\operatorname{int}W$, and
the same construction for the inverses, prove that this
is a joint continuous ordinary TOP loop $\Gamma$ on $B$.

Define $j_!$ precisely by the collapse map
\[
 B\longrightarrow W/\partial W,
\]
which is the quotient map on $W$ and sends its complement
to the basepoint. Equivalently it is the excision map
followed by forgetting support. The trace $T_\Gamma$
is identity off $W\times I$ and near its boundary.
Naturality of the relative universal lift difference
therefore gives
\[
          \eta_B[\Gamma]
          =\kappa_B(T_\Gamma)
          =j_!\kappa_W(T_\gamma)=j_!v.               \tag{L.7}
\]
At the cylinder level the cohomology map is
$H^3(N_W,\partial N_W)\to H^3(B\times I,B\times\partial I)$;
the interval suspension identifies it with the displayed
$j_!$ in degree two. The common identity collar makes
the two relative lifts literally identical where they
are glued, so no boundary class is lost.

The duality identity
\[
       \mathrm{PD}_B\circ j_!
                    =j_*\circ\mathrm{PD}_W           \tag{L.8}
\]
is the usual relative naturality for a codimension-zero
inclusion. It can also be seen directly in a finite PL
triangulation compatible with $W$ and its boundary:
extend a relative cochain by zero, and cap with the
fundamental chain of $B$. All simplices outside $W$
contribute zero; the remaining cap is the inclusion of
the cap with the relative fundamental chain of $W$.
This is over $\mathbb Z/2$, so no orientation sign enters.
Equations (L.6)--(L.8) prove (L.1).

If $j_*$ on $H_2$ is onto, (L.1) makes $\eta_B$ onto.
For every closed PL four-manifold $B$, A.136 identifies
the trace homomorphism with the composite
$\pi_1A_{\mathrm{TOP}}(B)\to\pi_0F_A
\xrightarrow{\pi_0\beta_B}H^2(B;\mathbb Z/2)$.
Thus its image lies in $J_B$, as defined in A.137.
A jointly PL loop has a jointly PL trace, so its
relative difference class is zero. These statements
require neither simple connectivity of $B$ nor the
assembly assumption in A.140. This gives
$J_B=H^2(B;\mathbb Z/2)$ and the stated automorphism
comparison conclusion. A.137 then has zero target for
$\lambda_B$, so its kernel formula gives
$\operatorname{im}(\pi_0c_P)=I_B$.

\textbf{Connected sums.}
For $B=C\#D$ as in the statement, take
$W=C\setminus\operatorname{int}D^4$ in the defining
connected-sum decomposition. It is a compact PL
submanifold with boundary $S^3$. Van Kampen applied
to $C=W\cup D^4$ shows $\pi_1(W)=\pi_1(C)=0$.
Removing the open four-ball does not change $H_2$.
The Mayer--Vietoris sequence for
\[
 B=(C\setminus\operatorname{int}D^4)
       \cup_{S^3}(D\setminus\operatorname{int}D^4)
\]
has zero adjacent groups $H_2(S^3;\mathbb Z/2)$ and
$H_1(S^3;\mathbb Z/2)$. Thus
$H_2(B;\mathbb Z/2)\cong H_2(W;\mathbb Z/2)
\oplus H_2(D;\mathbb Z/2)$, and the assumed vanishing
of the second summand makes $j_*$ an isomorphism.
This proves (L.2). Finally $S^1\times S^3$ has zero
second mod-$2$ homology, and the same Mayer--Vietoris
argument proves this for any finite connected sum of
such manifolds. The case $k=0$ is A.140 directly.
\end{proof}
